\section{Stability of the FSA-v2 SDE: derivations}\label{app:lyap}

This appendix supplies the derivations underlying the stability
analysis of Section~\ref{sec:lyap}. We compute the drift Jacobian
at the sedentary branch (Section~\ref{app:lyap-jacobian}), derive
the existence and local stability of the oscillatory branch
(Section~\ref{app:lyap-osc}), set out the bifurcation locus
$\Phi_c^{\mathrm{crit}}$ explicitly
(Section~\ref{app:lyap-bifurcation}), and prove the
Foster--Lyapunov drift inequality of Theorem~\ref{thm:lyap}
(Section~\ref{app:lyap-foster}).

\subsection{The drift Jacobian at the \texorpdfstring{$A=0$}{A=0} branch}\label{app:lyap-jacobian}

Write the deterministic drift vector field as
$f(B,F,A;\Phi_c) = (f_B,f_F,f_A)^{\!\top}$ with
\begin{align*}
  f_B(B,F,A) &= \kappa_B\,g_B(A)\,\Phi_c - B/\tau_B, \\
  f_F(B,F,A) &= \kappa_F\,\Phi_c - g_F(A)\,F/\tau_F, \\
  f_A(B,F,A) &= \mu(B,F)\,A - \eta\,A^3,
\end{align*}
where $g_B(A) = (1+\varepsilon_A A)/(1+\varepsilon_A A_{\mathrm{typ}})$
and $g_F(A) = (1+\lambda_A A)/(1+\lambda_A A_{\mathrm{typ}})$ are
the residual factors. Their derivatives at $A = 0$ are
$g_B'(0) = \varepsilon_A/(1+\varepsilon_A A_{\mathrm{typ}})$ and
$g_F'(0) = \lambda_A/(1+\lambda_A A_{\mathrm{typ}})$.

The Jacobian $J = \partial f/\partial(B,F,A)$ at $(B^*_0,F^*_0,0)$
has entries

\begin{itemize}
  \item Row $B$:
    $\partial f_B/\partial B = -1/\tau_B$,
    $\partial f_B/\partial F = 0$,
    $\partial f_B/\partial A = \kappa_B g_B'(0) \Phi_c
      = \kappa_B \varepsilon_A \Phi_c/(1+\varepsilon_A A_{\mathrm{typ}})$.
  \item Row $F$:
    $\partial f_F/\partial B = 0$,
    $\partial f_F/\partial F = -g_F(0)/\tau_F = -1/[\tau_F(1+\lambda_A A_{\mathrm{typ}})]$,
    $\partial f_F/\partial A = -g_F'(0) F^*_0/\tau_F
      = -\lambda_A F^*_0/[\tau_F(1+\lambda_A A_{\mathrm{typ}})]$.
  \item Row $A$:
    $\partial f_A/\partial B = \mu_B \cdot A = 0$ (at $A=0$),
    $\partial f_A/\partial F = (-\mu_F - 2\mu_{FF}(F-F_{\mathrm{typ}})) \cdot A = 0$ (at $A=0$),
    $\partial f_A/\partial A = \mu(B,F) - 3\eta A^2 = \mu(B^*_0,F^*_0)$ (at $A=0$).
\end{itemize}

In matrix form,
\begin{equation}
  J(B^*_0,F^*_0,0) \;=\;
    \begin{pmatrix}
      -1/\tau_B & 0 & \dfrac{\kappa_B \varepsilon_A \Phi_c}{1+\varepsilon_A A_{\mathrm{typ}}} \\[6pt]
      0 & -\dfrac{1}{\tau_F(1+\lambda_A A_{\mathrm{typ}})}
        & -\dfrac{\lambda_A F^*_0}{\tau_F(1+\lambda_A A_{\mathrm{typ}})} \\[6pt]
      0 & 0 & \mu(B^*_0,F^*_0)
    \end{pmatrix}.
  \label{eq:jacobian}
\end{equation}
This is upper triangular, so its eigenvalues are the diagonal
entries:
\begin{equation}
  \lambda_1 = -1/\tau_B,
  \qquad
  \lambda_2 = -1/[\tau_F(1+\lambda_A A_{\mathrm{typ}})],
  \qquad
  \lambda_3 = \mu(B^*_0,F^*_0).
  \label{eq:eigs}
\end{equation}
The first two are unconditionally negative (with $\tau_B,\tau_F > 0$
and $\lambda_A,A_{\mathrm{typ}} \ge 0$); the third has the sign of
the Stuart--Landau growth rate. By Lyapunov's first method, the
$A=0$ equilibrium is locally asymptotically stable iff all three
eigenvalues are negative, i.e.\ iff $\mu(B^*_0,F^*_0) < 0$. This
proves Proposition~\ref{prop:stab-zero}. \qed

\subsection{The oscillatory branch: existence and local stability}\label{app:lyap-osc}

The oscillatory equilibrium of the deterministic skeleton sits at
$A^*_{\mathrm{osc}} = \sqrt{\mu(B^*,F^*)/\eta} > 0$ with
$(B^*,F^*)$ given implicitly by setting
$\dot B = \dot F = 0$ in
\eqref{eq:fsa-B}--\eqref{eq:fsa-F} at the equilibrium amplitude
$A = A^*_{\mathrm{osc}}$. The implicit dependence on $A$ enters
through the residual factors $g_B(A),g_F(A)$ of
Section~\ref{app:lyap-jacobian}.

To leading order, and consistent with the
operating-point linearisation of Definition~\ref{def:operating-point},
we evaluate the residual factors at $A = A_{\mathrm{typ}}$, where
$g_B(A_{\mathrm{typ}}) = g_F(A_{\mathrm{typ}}) = 1$. Then
$(B^*,F^*) \approx (B^*_0,F^*_0)$ as in \eqref{eq:eq-zero}, and
\begin{equation}
  A^*_{\mathrm{osc}}(\Phi_c)
   \;\approx\; \sqrt{\,\mu(B^*_0(\Phi_c),F^*_0(\Phi_c))\,/\,\eta\,},
  \label{eq:A-osc-approx}
\end{equation}
defined wherever $\mu(B^*_0,F^*_0) > 0$. This is the formula used
to populate Table~\ref{tab:goldilocks} of the main text and the
right-hand panel of Figure~\ref{fig:goldilocks}.

\paragraph{Local stability.} A short Jacobian computation -- entirely
analogous to that of Section~\ref{app:lyap-jacobian} but evaluated at
$(B^*,F^*,A^*_{\mathrm{osc}})$ -- shows that the eigenvalue along
the $A$-direction at the oscillatory branch is
\begin{equation*}
  \partial_A\bigl[\mu(B,F)\,A - \eta A^3\bigr]\Big|_{A=A^*_{\mathrm{osc}}}
  \;=\; \mu(B^*,F^*) - 3\eta\,(A^*_{\mathrm{osc}})^2
  \;=\; \mu(B^*,F^*) - 3\mu(B^*,F^*) \;=\; -2\,\mu(B^*,F^*).
\end{equation*}
Since the oscillatory branch exists only when $\mu(B^*,F^*) > 0$,
the $A$-direction eigenvalue is strictly negative there. The other
two eigenvalues remain $-1/\tau_B$ and
$-1/[\tau_F(1+\lambda_A A_{\mathrm{typ}})]$ (with corrections
$O(\varepsilon_A,\lambda_A)$ from the residual factors at
$A^*_{\mathrm{osc}} \neq A_{\mathrm{typ}}$). All three eigenvalues
are negative, so the oscillatory branch is locally asymptotically
stable wherever it exists.

\paragraph{Maximum of $A^*_{\mathrm{osc}}$.} Differentiating
\eqref{eq:A-osc-approx} with respect to $\Phi_c$ and using
$B^*_0(\Phi_c)$ saturating at unity for $\Phi_c \ge 1/\tilde b
\approx 1.98$,
\begin{equation*}
  \frac{\dd \mu}{\dd \Phi_c}
   = \begin{cases}
     \mu_B\,\tilde b - \mu_F\,\tilde f - 2\mu_{FF}\,\tilde f\,(\tilde f\Phi_c - F_{\mathrm{typ}})
        & 0 \le \Phi_c < 1/\tilde b \\[4pt]
     -\mu_F\,\tilde f - 2\mu_{FF}\,\tilde f\,(\tilde f\Phi_c - F_{\mathrm{typ}})
        & \Phi_c \ge 1/\tilde b.
   \end{cases}
\end{equation*}
For the truth values of Table~\ref{tab:fsa-truth}, the first branch
gives a positive derivative throughout $[0, 1/\tilde b)$ -- i.e.\
$\mu$ is monotone increasing in this range -- and the second
branch is negative for all $\Phi_c \ge 1/\tilde b$. The maximum
therefore occurs exactly at the saturation point
\begin{equation}
  \Phi_{\mathrm{opt}} \;=\; 1/\tilde b
   \;=\; \frac{1+\varepsilon_A A_{\mathrm{typ}}}{\kappa_B \tau_B}
   \;\approx\; 1.99,
  \label{eq:Phi-opt}
\end{equation}
giving $\mu_{\max} \approx 0.207$ and
$A^*_{\mathrm{osc}}(\Phi_{\mathrm{opt}}) \approx 1.02$.

\subsection{The bifurcation locus}\label{app:lyap-bifurcation}

Above $\Phi_{\mathrm{opt}}$, $B$ remains saturated at unity and
$\mu$ depends only on $F = \tilde f \Phi_c$:
\begin{equation}
  \mu(\Phi_c) \;=\; \mu_0 + \mu_B - \mu_F\,\tilde f\,\Phi_c - \mu_{FF}(\tilde f\Phi_c - F_{\mathrm{typ}})^2,
  \quad \Phi_c \ge \Phi_{\mathrm{opt}}.
  \label{eq:mu-of-Phi}
\end{equation}
This is a downward-opening quadratic in $\Phi_c$ with leading
coefficient $-\mu_{FF}\tilde f^2$. At the truth values of
Table~\ref{tab:fsa-truth},
\begin{equation*}
  \tilde b \;=\; \frac{0.01248 \cdot 42}{1.04} \;\approx\; 0.504,
  \qquad
  \tilde f \;=\; 0.030 \cdot 6.364 \cdot 1.1 \;\approx\; 0.210.
\end{equation*}
Setting \eqref{eq:mu-of-Phi} to zero and solving yields a unique
positive root in the saturated regime:
\begin{equation}
  \Phi_c^{\mathrm{crit}} \;\approx\; 3.71\quad\text{(TRIMP-units per day)}.
\end{equation}
For $\Phi_{\mathrm{opt}} < \Phi_c < \Phi_c^{\mathrm{crit}}$ we have
$\mu > 0$ and the oscillatory branch persists (with decreasing
amplitude); for $\Phi_c > \Phi_c^{\mathrm{crit}}$ we have $\mu < 0$
and the oscillatory branch ceases to exist, leaving the sedentary
branch as the unique stable equilibrium. The bifurcation at
$\Phi_c^{\mathrm{crit}}$ is a Hopf-like merger of the oscillatory
and sedentary branches.

In the under-trained regime $0 \le \Phi_c < \Phi_{\mathrm{opt}}$,
$\mu(B^*_0,F^*_0)$ is given by the unsaturated formula
$\mu_0 + \mu_B \tilde b\,\Phi_c - \mu_F\,\tilde f\,\Phi_c
       - \mu_{FF}(\tilde f\Phi_c - F_{\mathrm{typ}})^2$
which we have already used to compute
Table~\ref{tab:goldilocks}; this remains positive throughout
$[0,\Phi_{\mathrm{opt}}]$, so the oscillatory branch exists and is
stable for every constant load below the bifurcation threshold.

\subsection{Foster--Lyapunov drift inequality (proof of Theorem~\ref{thm:lyap})}\label{app:lyap-foster}

We now prove that the quadratic Lyapunov function
$V(B,F,A) = \tfrac{1}{2}(a B^2 + b F^2 + c A^2)$ of
\eqref{eq:V-quadratic} satisfies the Foster--Lyapunov drift
inequality \eqref{eq:foster-lyap} on the operating domain
$\Xcal = [0,1] \times [0,\infty)^2$, with rate
$\alpha = \min\!\bigl(2/\tau_B,\,1/[\tau_F(1+\lambda_A A_{\mathrm{typ}})]\bigr)$
and a finite constant $C > 0$.

Let $\Lcal$ denote the infinitesimal generator of the SDE: for any
$V \in C^2(\Xcal)$,
\begin{equation}
  \Lcal V(x) \;=\;
    \sum_{i \in \{B,F,A\}} b_i(x)\,\partial_i V(x)
    \;+\; \tfrac{1}{2}\sum_i a_i(x)\,\partial_{ii}^2 V(x),
  \label{eq:generator}
\end{equation}
with $b_i$ the $i$th drift component and
$a_B(x) = \sigma_B^2 B(1-B)$,
$a_F(x) = \sigma_F^2 F$,
$a_A(x) = \sigma_A^2 A$. We bound the three components separately.
Throughout, $K_\bullet$ denotes a finite constant depending only on
the truth parameters, $\Phimax$, and the choice of $a,b,c$.

\paragraph{$B$-component.} Using $B \in [0,1]$ and the fact that
$g_B(A) = (1+\varepsilon_A A)/(1+\varepsilon_A A_{\mathrm{typ}})$
is bounded above on the operating domain by some
$g_{B,\max} < \infty$ (depending on the practical upper bound on
$A$ that we wish to enforce; on the operating domain
$A \in [0,A_{\max}]$ with any chosen $A_{\max}$, set
$g_{B,\max} = (1+\varepsilon_A A_{\max})/(1+\varepsilon_A A_{\mathrm{typ}})$),
\begin{align*}
  a B \cdot \bigl[\kappa_B g_B(A) \Phi - B/\tau_B\bigr]
   + \tfrac{a}{2} \sigma_B^2 B(1-B)
  &\le \;a\,\kappa_B\,g_{B,\max}\,\Phimax\,B - (a/\tau_B)\,B^2 + \tfrac{a}{8}\sigma_B^2 \\
  &\le \;a\,\kappa_B\,g_{B,\max}\,\Phimax + \tfrac{a}{8}\sigma_B^2 - (a/\tau_B)\,B^2,
\end{align*}
where the second line uses $B \le 1$ and the bound
$B(1-B) \le 1/4$ on $[0,1]$. Define
$K_B := a\kappa_B g_{B,\max}\Phimax + a\sigma_B^2/8$.

\paragraph{$F$-component.} Using $g_F(A) \ge g_F(0) = 1/(1+\lambda_A A_{\mathrm{typ}})$
for $A \ge 0$, and writing $\tau_F^{\mathrm{old}} := \tau_F(1+\lambda_A A_{\mathrm{typ}})$,
\begin{align*}
  bF \cdot \bigl[\kappa_F\Phi - g_F(A)F/\tau_F\bigr]
   + \tfrac{b}{2}\sigma_F^2 F
  &\le \;b\,\kappa_F\,\Phimax\,F - bF^2/\tau_F^{\mathrm{old}} + \tfrac{b}{2}\sigma_F^2 F \\
  &= \; b F\,(\kappa_F\Phimax + \sigma_F^2/2) - bF^2/\tau_F^{\mathrm{old}}.
\end{align*}
Apply Young's inequality $xy \le x^2/(2\beta) + \beta y^2/2$ to the
linear-in-$F$ term, with $\beta = \tau_F^{\mathrm{old}}$ and
$y = (\kappa_F\Phimax + \sigma_F^2/2)$:
\[
  b F (\kappa_F\Phimax + \sigma_F^2/2)
   \;\le\; bF^2/(2\tau_F^{\mathrm{old}})
           + (b\tau_F^{\mathrm{old}}/2)(\kappa_F\Phimax + \sigma_F^2/2)^2.
\]
Combined,
\[
  bF \cdot \bigl[\kappa_F\Phi - g_F(A)F/\tau_F\bigr] + \tfrac{b}{2}\sigma_F^2 F
   \;\le\; -\bigl(b/(2\tau_F^{\mathrm{old}})\bigr)F^2 + K_F,
\]
with $K_F := (b\tau_F^{\mathrm{old}}/2)(\kappa_F\Phimax + \sigma_F^2/2)^2$.

\paragraph{$A$-component.} The growth rate satisfies
$\mu(B,F) \le \mu_{\max} := \mu_0 + \mu_B$ for $B \in [0,1]$ and
$F \ge 0$, since the contributions $-\mu_F F$ and
$-\mu_{FF}(F-F_{\mathrm{typ}})^2$ are non-positive. Therefore
\begin{equation*}
  cA\bigl[\mu(B,F)\,A - \eta A^3\bigr] + \tfrac{c}{2}\sigma_A^2 A
   \;\le\; -c\eta A^4 + c\mu_{\max} A^2 + \tfrac{c}{2}\sigma_A^2 A.
\end{equation*}
Two applications of Young's inequality remove the unwanted $A^2$
and linear-$A$ terms:

\emph{(i)} For the $A^4$-vs-$A^2$ pair, complete the square with
$x = A^2$:
\[
  -c\eta A^4 + c\mu_{\max} A^2
   = -c\eta\bigl(A^2 - \mu_{\max}/(2\eta)\bigr)^2 + c\mu_{\max}^2/(4\eta)
   \;\le\; -\tfrac{c\eta}{2}A^4 + \tfrac{c\mu_{\max}^2}{2\eta},
\]
where we discarded half of the negative quartic contribution to
absorb the maximum of the non-negative quadratic.

\emph{(ii)} For the linear $A$-term against the diffusion: by the
elementary $A \le (\eta/(2\sigma_A^2))A^4 + \sigma_A^2/(2\eta)$
(setting $\alpha = \eta/(2\sigma_A^2)$ in
$x \le \alpha x^4 + 1/(4\alpha)$ -- proved by completing the
square on $x \mapsto \alpha x^4 - x$ and verifying the minimum is
$-1/(4\alpha)$, so $x - \alpha x^4 \le 1/(4\alpha)$, i.e.\
$x \le \alpha x^4 + 1/(4\alpha)$), and we then absorb the $A^4$
piece into a small fraction of the $-\tfrac{c\eta}{2}A^4$ term:
\[
  \tfrac{c}{2}\sigma_A^2 A
   \;\le\; \tfrac{c}{2}\sigma_A^2 \cdot
            \bigl(\eta/(2\sigma_A^2) A^4 + \sigma_A^2/(2\eta)\bigr)
   \;=\; \tfrac{c\eta}{4} A^4 + \tfrac{c\sigma_A^4}{4\eta}.
\]
Wait -- that produces the wrong sign on $A^4$. We instead want the
linear-$A$ term to sit on top of the cubic damping; redo with the
opposite split. Use $A \le 1 + A^2$ (a strict bound for $A \ge 0$
verified by completing the square on $1 + A^2 - A$); then
\[
  \tfrac{c}{2}\sigma_A^2 A
   \;\le\; \tfrac{c\sigma_A^2}{2}(1 + A^2)
   \;=\; \tfrac{c\sigma_A^2}{2} + \tfrac{c\sigma_A^2}{2}A^2,
\]
and the residual $A^2$ is then absorbed via the same complete-the-square
argument as in step~(i), but with $\mu_{\max}$ replaced by
$\sigma_A^2/2$. Combining all three pieces,
\[
  cA[\mu(B,F)A - \eta A^3] + \tfrac{c}{2}\sigma_A^2 A
   \;\le\; -\tfrac{c\eta}{4}A^4
            + \tfrac{c\mu_{\max}^2}{2\eta}
            + \tfrac{c\sigma_A^2}{2}
            + \tfrac{c\sigma_A^4}{8\eta}
   \;=:\; -\tfrac{c\eta}{4}A^4 + K_A.
\]

\paragraph{Combining the three components.} Adding the three bounds,
\begin{equation*}
  \Lcal V(x)
   \;\le\;
   -(a/\tau_B)\,B^2
   - \bigl(b/(2\tau_F^{\mathrm{old}})\bigr) F^2
   - (c\eta/4)\,A^4
   + (K_B + K_F + K_A).
\end{equation*}
The first two negative terms already match $-\alpha V$ with
$\alpha = \min(2/\tau_B,\,1/\tau_F^{\mathrm{old}})$. For the third
we use the elementary inequality
\[
  -(c\eta/4)\,A^4 \;\le\; -(c\alpha/2)\,A^2 + c\alpha^2/(4\eta),
\]
which follows by completing the square in $A^2$ on the parabola
$(c\eta/4)x^2 - (c\alpha/2)x$. Plugging this in,
\begin{equation*}
  \Lcal V(x)
   \;\le\;
   -\tfrac{\alpha}{2}\bigl(a B^2 + b F^2 + c A^2\bigr)
   + \bigl(K_B + K_F + K_A + c\alpha^2/(4\eta)\bigr)
   \;=\;
   -\alpha\,V(x) + C,
\end{equation*}
where $C := K_B + K_F + K_A + c\alpha^2/(4\eta)$ is finite. This is
the required drift inequality \eqref{eq:foster-lyap}, completing
the proof of Theorem~\ref{thm:lyap}. \qed

\medskip

The exponential ergodicity statement of Theorem~\ref{thm:ergodic}
then follows from the standard Foster--Lyapunov $\Rightarrow$ V-uniform
ergodicity machinery in Khasminskii (1980) or, in discrete-time
language transferable to continuous time via the resolvent,
Meyn--Tweedie (1993, Theorem~6.1). The constants $K_0$ and $\rho$
in \eqref{eq:ergodic-bound} can be made explicit at the cost of
further notation; we do not pursue this here.
